% =============================================================================
% Multi-Asset Statistical Arbitrage via Regularized Regression, Nonlinear
% Extensions, and a Low-Latency C++ Execution Engine
%
% CANONICAL SOURCE OF TRUTH for the report. Built by report/build_report.py,
% which (a) regenerates figures/tables from results/, (b) compiles this to PDF
% when a TeX engine is available, and (c) always renders an HTML preview.
%
% Compile anywhere with a normal TeX distribution:
%   latexmk -pdf report.tex     (or: pdflatex -> biber -> pdflatex x2)
% =============================================================================
\documentclass[11pt]{article}

% --- Packages (per spec Section 6.2) ---
\usepackage[margin=1in]{geometry}
\usepackage[T1]{fontenc}
\usepackage{lmodern}
\usepackage{amsmath,amssymb,amsthm}
\usepackage{booktabs}
\usepackage{graphicx}
\usepackage{subcaption}
\usepackage[table]{xcolor}
\usepackage{multirow}
\usepackage{enumitem}
\usepackage{float}
\usepackage{threeparttable}
\usepackage[colorlinks=true,linkcolor=blue,citecolor=blue,urlcolor=blue]{hyperref}
\usepackage[style=authoryear,maxcitenames=2,sorting=nyt,backend=biber]{biblatex}
\addbibresource{references.bib}

% --- Theorems/environments ---
\newtheorem{assumption}{Assumption}[section]
\newtheorem{proposition}{Proposition}[section]

% --- Convenience commands ---
\newcommand{\E}{\mathbb{E}}
\newcommand{\V}{\mathbb{V}}
\newcommand{\R}{\mathbb{R}}
\newcommand{\eps}{\varepsilon}
\newcommand{\half}{\tfrac{1}{2}}
\DeclareMathOperator{\Var}{Var}
\DeclareMathOperator{\Cov}{Cov}

% --- Title ---
\title{\bfseries Multi-Asset Statistical Arbitrage via Regularized Regression,\\
       Nonlinear Extensions, and a Low-Latency C++ Execution Engine}
\author{Research and Engineering Deliverable\\[2pt] \small Autonomous Agent Project}
\date{\today}

\begin{document}
\maketitle

% =============================================================================
% 1. ABSTRACT
% =============================================================================
\section*{Abstract}
\label{sec:abstract}
\emph{[Filled in the final Part. One paragraph: the question, the method, and the
headline -- honest -- finding. States whether, after multiple-testing correction,
transaction costs, and out-of-sample validation, a statistically defensible
statistical-arbitrage edge exists in the studied universe, and under what conditions.]
}
\vspace{2em}

\tableofcontents
\newpage

% =============================================================================
% 2. INTRODUCTION & MOTIVATION
% =============================================================================
\section{Introduction and Motivation}
\label{sec:intro}

Economic rationale for basket statistical arbitrage. A tradeable residual is defined
relative to an economically motivated basket of related assets, not an arbitrary
regressor set. Position: short the target / long the basket (or the reverse) when the
residual deviates beyond a statistically justified threshold; the position is sized by
the estimated weights.

\textbf{Literature context.}
\begin{itemize}
  \item \cite{avellaneda2010statistical}: systematic factor-based stat-arb with PCA
    and sector ETFs; the central reference for basket selection and residual dynamics.
  \item \cite{dofaff2010does}: does pairs trading still work after costs; emphasizes
    cost sensitivity and mean-reversion half-life.
  \item \cite{gatev2006pairs}: the canonical pairs-trading profitability study and its
    critique of data-snooping (top-decile selection bias).
  \item \cite{dieboldmariano1995}: forecast comparison methodology used here.
  \item \cite{white2000realitycheck}, \cite{hansen2005test}: data-snooping tests
    (White Reality Check, Hansen's SPA) used to answer ``is the best Sharpe real?''
\end{itemize}

Our contribution: a walk-forward, cost-aware, C++-accelerated, adversarial
self-critiqued version of the basket stat-arb idea, with explicit multiple-testing
correction and honest discussion of what it does and does not replicate.

\subsection{The base model}
\label{sec:basemodel}
For a target asset $Y_t$ regressed on a basket of $n$ candidate assets $X_{1,t},\dots,X_{n,t}$:
\begin{equation}
  Y_t = \alpha + \sum_{i=1}^{n}\beta_i\,X_{i,t} + \eps_t ,
  \label{eq:basemodel}
\end{equation}
where $\eps_t$ is the tradeable signal. The specification of which $X$'s to include,
how to estimate $\beta$, and how to map $\eps_t$ to positions is the substance of this
report.

% =============================================================================
% 3. DATA
% =============================================================================
\section{Data}
\label{sec:data}

\subsection{Source and acquisition}
\label{sec:data:source}
The in-sandbox primary dataset is the classic \emph{NYSE (Kaggle)} dataset,
obtained from the GitHub repository
\href{https://github.com/ashishpatel26/NYSE-STOCK_MARKET-ANALYSIS-USING-LSTM}{\texttt{ashishpatel26/NYSE-STOCK\_MARKET-ANALYSIS-USING-LSTM}}
(downloaded as a repository tarball via the GitHub API on \texttt{\today}).
It contains daily price data for 501 U.S. securities (an S\&P 500-style snapshot),
with fields \texttt{date, symbol, open, close, low, high, volume}, plus a
separate split-adjusted price file, a \texttt{securities} file with GICS sector /
sub-industry tags, and a quarterly \texttt{fundamentals} file.

\emph{Environment note.} This sandbox cannot reach external market-data endpoints
(Yahoo Finance, stooq, Nasdaq) at runtime; GitHub is the only reachable data host.
A companion downloader
(\texttt{scripts/data/download.py} with \texttt{--source yahoo}) is provided so a
$\geq\!10$-year panel (2010--present) can be produced on a normal internet
connection; the cleaning and QA pipeline below is identical for both sources. The
remainder of this section reports on the in-sandbox GitHub dataset actually used for
the analyses in this report.

\subsection{Cleaning and adjustment methodology}
\label{sec:data:cleaning}
All raw files are ingested by \texttt{scripts/data/clean.py} into a uniform long
panel with columns
\[
\texttt{date},\ \texttt{ticker},\ \texttt{open},\ \texttt{high},\ \texttt{low},\ \texttt{close},\ \texttt{adj\_close},\ \texttt{volume}.
\]
Corporate-action handling: for the NYSE source we use the provided split-adjusted
prices; these are \emph{split}-adjusted only, so all returns in this report are
\emph{price} (split-adjusted) returns, not total returns -- a documented
limitation. (When the Yahoo path is used, \texttt{adj\_close} is fully adjusted for
splits \emph{and} dividends.) Rows with non-positive prices are dropped, and the
negligible fraction of zero-volume rows (0.002\% in the source) is removed. Duplicate
(date, ticker) pairs are dropped (last value wins) and the panel is sorted.

\subsection{Universe description}
\label{sec:data:universe}
The universe is the 501-name S\&P 500-style snapshot, mapped to 11 GICS sectors.
A data-quality gate (\texttt{scripts/data/qa.py}) verifies positive prices, no
missing closes, no duplicate dates, sane volumes, per-symbol span, and calendar-gap
tolerance; it passes 10/10 checks. Summary statistics are regenerated into
\autoref{tab:data_summary} directly from the cleaned panel.

\input{tables/data_section.tex}

\subsection{Survivorship-bias discussion}
\label{sec:data:survivorship}
This dataset is a 2016-era constituent snapshot with a $\approx\!7$-year daily span
(2010--2016). Two biases must be disclosed. First, because the panel was assembled
from a 2016 membership list, names delisted or acquired before 2016 are absent;
second, the panel is short of the $10$-year minimum span in the specification. The
constituent \texttt{securities} file records each name's \emph{date first added}, which
partially identifies membership changes, but a fully survivorship-bias-free panel
requires the Yahoo-path download (10+ years) plus a historical constituent list
(e.g.\ the Wikipedia/BlackRock membership history), which we list as the recommended
upgrade. All results in later sections must therefore be read as conditional on this
universe and corrected for the multiple tests it entails.

% =============================================================================
% 4. METHODOLOGY
% =============================================================================
\section{Methodology}
\label{sec:method}

\subsection{Estimator derivations}
\label{sec:estimators}
\emph{[Full derivations filled in Part 2/5.]}

\subsubsection{Ordinary least squares}
Closed-form estimator
\begin{equation}
  \hat\beta = \bigl(X^{\top}X\bigr)^{-1}X^{\top}Y ,
  \label{eq:ols}
\end{equation}
derived by minimizing the residual sum of squares. Gauss--Markov assumptions and what
happens when they are violated in financial time series (non-i.i.d. errors,
heteroskedasticity, serial correlation) are treated here.

\subsubsection{Ridge regression}
\begin{equation}
  \hat\beta_{\mathrm{ridge}} = \bigl(X^{\top}X + \lambda I\bigr)^{-1}X^{\top}Y ,
  \label{eq:ridge}
\end{equation}
with geometric and Bayesian (MAP under a Gaussian prior) interpretations.

\subsubsection{Lasso}
Via subgradient conditions; no closed form. Coordinate descent / LARS is described.

\subsubsection{Principal component regression}
Via the eigendecomposition of the covariance matrix; connection to market/sector
factor structure.

\subsection{Statistical diagnostics}
\label{sec:diagnostics}
\emph{[Definitions, null hypotheses, and joint interpretation filled in Part 2/6.]}
Stationarity (ADF, KPSS), cointegration (Engle--Granger, Johansen), residual
autocorrelation (Durbin--Watson, Ljung--Box), heteroskedasticity (Breusch--Pagan /
White), Newey--West HAC correction, and the Ornstein--Uhlenbeck mean-reversion fit
with half-life $\ln(2)/\theta$ compared to the holding period.

\subsection{Multiple-testing correction}
\label{sec:multtest}
Bonferroni, White Reality Check, Hansen's SPA. \emph{[Numeric result in Part 10.]}

\subsection{Forecast comparison}
Diebold--Mariano test across model pairs.

% =============================================================================
% 5. MODEL COMPARISON RESULTS
% =============================================================================
\section{Model Comparison Results}
\label{sec:modelcomparison}
OLS / Ridge / Lasso / Elastic Net / PCA / Kalman. In-sample and out-of-sample $R^2$,
coefficient stability over time, computational cost. \emph{[Filled in Parts 5, 8.]}

% =============================================================================
% 6. NONLINEAR & ADAPTIVE EXTENSIONS
% =============================================================================
\section{Nonlinear and Adaptive Extensions}
\label{sec:nonlinear}
Gradient-boosted trees / shallow network, RESET test, learning curves, honest
bias--variance assessment; Kalman-filtered time-varying beta. \emph{[Parts 8, 9.]}

% =============================================================================
% 7. BACKTEST RESULTS
% =============================================================================
\section{Backtest Results}
\label{sec:backtest}
Gross vs.\ net of costs, cost sensitivity, regime breakdown, bootstrap confidence
intervals. \emph{[Parts 4, 7.]}

% =============================================================================
% 8. ROBUSTNESS AND STATISTICAL VALIDITY
% =============================================================================
\section{Robustness and Statistical Validity}
\label{sec:robustness}
Multiple-testing correction numbers, Diebold--Mariano results, sensitivity heatmaps.
\emph{[Parts 6, 10.]}

% =============================================================================
% 9. SYSTEM ARCHITECTURE & LATENCY ANALYSIS
% =============================================================================
\section{System Architecture and Latency Analysis}
\label{sec:architecture}
C++ design rationale, cache-efficiency benchmarks, latency histograms.
\emph{[Parts 3, 11.]}

% =============================================================================
% 10. DISCUSSION: WHAT THIS DOES AND DOESN'T REPLICATE
% =============================================================================
\section{Discussion: What This Does and Does Not Replicate}
\label{sec:nonreplicate}
\emph{[Filled in Part 12.] Explicit, specific: no proprietary alternative data; no
colocated / FPGA-level execution; no cross-asset-class integration; no live capital
allocation feedback loops -- and why each matters at real fund level.]}

% =============================================================================
% 11. CONCLUSION
% =============================================================================
\section{Conclusion}
\label{sec:conclusion}
\emph{[Filled in Part 12.] Clear, honest statement of whether a statistically defensible
edge exists after all corrections and costs, and under what conditions.]}

% =============================================================================
% 12. REFERENCES
% =============================================================================
\printbibliography

% =============================================================================
% 13. APPENDIX
% =============================================================================
\appendix
\section{Additional Derivations and Diagnostic Tables}
\label{sec:appendix}
\emph{[Extended derivations and extra diagnostic tables in later Parts.]}

\end{document}
